\documentclass{article}
\usepackage{amsfonts,amsthm,amsmath}
\usepackage[mathscr]{euscript}
\usepackage{enumitem}

\setcounter{section}{0}
\renewcommand{\thesection}{\S\arabic{section}}
\renewcommand{\thesubsection}{\arabic{section}}

\newtheorem{thm}{Theorem}
\newenvironment{theorem}[1]{
	\renewcommand{\thethm}{#1}
	\begin{thm}
		}{\end{thm}}

\newtheorem{lma}{Lemma}
\newenvironment{lemma}[1]{
	\renewcommand{\thelma}{#1}
	\begin{lma}
		}{\end{lma}}

\theoremstyle{definition}
\newtheorem{ex}{Exercise}
\newenvironment{exercise}[1]{
	\renewcommand{\theex}{#1}
	\begin{ex}
		}{\end{ex}}

\title{Exercises 8.C}
\author{Connor Mooneyhan}
\date{February 28, 2024}

\begin{document}

\maketitle

\begin{ex}
	Suppose $T \in \mathcal{L}(\mathbb{C}^4)$ is such that the eigenvalues of $T$ are $3, 5, 8$.
	Prove that $(T - 3I)^2(T - 5I)^2(T - 8I)^2 = 0$.
\end{ex}
\begin{proof}
	Since $\mathrm{dim}\,\mathbb{C}^4 = 4$, the degree of the minimal polynomial $p$ of $T$ is at most $4$.
	The degree of the above polynomial, call it $q$, is $6$.
	Since $p(T) = 0$, the zeroes of $p$ are exactly the eigenvalues of $T$, namely $3, 5, 8$.
	This means that $p$ can be written as \begin{equation*}
		p(z) = (z - 3)^{n_1}(z - 5)^{n_2}(z - 8)^{n_3}
	\end{equation*} where $n_1, n_2, n_3 \in \lbrace 1, 2 \rbrace$.
	No matter what the values of those powers end up being, each is at most $2$, so $q$ is clearly a polynomial multiple of $p$.
	By 8.46, this means $q(T) = 0$.
\end{proof}

\begin{ex}
	Suppose $V$ is a complex vector space. Suppose $T \in \mathcal{L}(V)$ is such that $5$ and $6$ are eigenvalues of $T$ and that $T$ has no other eigenvalues. Prove that $(T - 5I)^{n - 1}(T - 6I)^{n - 1} = 0$, where $n = \mathrm{dim}\,V$.
\end{ex}
\begin{proof}
	We know that the minimal polynomial, $p$, of $T$ is of degree at most $\mathrm{dim}\,V$.
	We can also see that given that $T$ has two distinct eigenvalues, $\mathrm{dim}\,V \geq 2$, and the multiplicity of each eigenvalue is at most $n - 1$ since their multiplicities must each be at least $1$ and their sum must equal $n$ because $V$ is a complex vector space.
	Thus, the characteristic polynomial $q$ can be described as \begin{equation*}
		q(z) = (z - 5)^{m_1}(z - 6)^{m_2}
	\end{equation*} where $m_1, m_2 \in \lbrace 1, \dots, n - 1 \rbrace$.
	Given this, clearly the polynomial in question (call it $r$) is a polynomial multiple of $q$, and since $q(T) = 0$, so is $r$.
\end{proof}

\begin{ex}
	Give an example of an operator on $\mathbb{C}^4$ whose characteristic polynomial equals $(z - 7)^2(z - 8)^2$.
\end{ex}
\begin{proof}
	The operator $T \in \mathcal{L}(V)$ defined by \begin{equation*}
		T(z_1, z_2, z_3, z_4) = (7z_1, 7z_2, 8z_3, 8z_4)
	\end{equation*} clearly has as its eigenvalues $7$ and $8$, with the corresponding eigenspaces each being two-dimensional.
	Since that makes the sum of the dimension of the eigenspaces $4$ (which is $\mathrm{dim}\,V$), the dimension of the generalized eigenspaces must also each be $2$, and thus the multiplicity of each eigenvalue is $2$.
	It follows by definition that the characteristic polynomial of $T$ is $(z - 7)^2(z -8)^2$ as desired.
\end{proof}

\setcounter{ex}{7}

\begin{ex}
	Suppose $T \in \mathcal{L}(V)$. Prove that $T$ is invertible if and only if the constant term in the minimal polynomial of $T$ is nonzero.
\end{ex}
\begin{proof}
	Denote the minimal polynomial of $T$ as $p(z) = a_0 + a_1z + \dots + a_nz^n$.
	First, assume that $a_0 \neq 0$.
	Then \begin{align*}
		a_0      I & = -a_1T + \dots + a_nT^n                                   \\
		\implies I & = -\frac{a_1}{a_0}T + \dots + \frac{a_n}{a_0}T^n           \\
		           & = T(-\frac{a_1}{a_0}I + \dots + \frac{a_n}{a_0}T^{n - 1}),
	\end{align*} so $T$ is invertible.

	Now assume that $a_0 = 0$. Then we have \begin{align}
		0 & = a_1T + \dots + a_nT^n           \\
		  & = T(a_1I + \dots + a_nT^{n - 1}).
	\end{align}
  For convenience, we define $q \in \mathcal{P}(\mathbf{F})$ by $q(z) = a_1 + \dots + a_nz^{n - 1}$ so that the right side of (2) becomes $T(q(T))$.
	Since $q$ is a polynomial of degree less than $n $, $q$ cannot be the zero map.
  Thus there exists some $v \in V$ for which $(q(T))v \neq 0$.
  Equation (2) indicates that $(q(T))v \in \mathrm{null}\,T$, so $T$ is not invertible.
\end{proof}

\setcounter{ex}{11}

\begin{ex}
  Suppose $V$ is a complex vector space and $T \in \mathcal{L}(V)$. Prove that $V$ has a basis consisting of eigenvectors of $T$ if and only if the minimal polynomial of $T$ has no repeated zeroes.
\end{ex}
\begin{proof}
  We begin with the easy direction.
  If $V$ has a basis consisting of eigenvectors of $T$, then given eigenvalues $\lambda_0, \dots, \lambda_n$, the polynomial \begin{equation*}
    p(T) = (T-\lambda_0I)\dots(T-\lambda_nI)
  \end{equation*} sends every basis element to zero, and thus is the zero map.
  $p$ must then be the minimal polynomial because taking away any factors would take away zeroes.

  For the other direction, assume that the minimal polynomial of $T$ has no repeated zeroes.
  Then write the minimal polynomial $p$ as \begin{equation*}
    p(T) = (T-\lambda_0)\dots(T-\lambda_n)
  \end{equation*} where each $\lambda_j$ is a distinct zero of $p$.
  Since $p$ is the minimal polynomial of $T$, $p$ sends each basis element to 0.
\end{proof}

\end{document}
